\subsection{局部化数系异常相位的 pairwise-union 分解与 prime-register 重数公式}\label{subsec:conclusion-localized-anomaly-pairwise-union-prime-ledger}
在小节 \ref{subsec:conclusion-localized-groups-mayer-vietoris-end-kernel} 与小节 \ref{subsec:conclusion-localized-torsion-godel-fixedpoint-cantor-geometry} 已经把局部化数系 \(G_S=\ZZ[S^{-1}]\) 的加法拼合、自同构核公式与对偶 torsion 轮廓固定之后，再把相位异常接口限制到有限个 rank-\(1\) 局部化轴上，还可得到一条更强的刚性链：单一局部化轴不产生任何连续二阶异常；有限个局部化轴的全部连续异常恰由 pairwise unions 生成；三寄存器情形不出现新的 primitive 三体连续异常；而相应 prime-register 核的逐素数重数可用一条显式二项式公式完全写出。

\begin{theorem}[单一局部化数系的连续异常相位严格消失]\label{thm:conclusion-localized-single-axis-anomaly-vanishing}
\leanverified{paper\_conclusion\_localized\_single\_axis\_anomaly\_vanishing}
设 \(S\subset\mathbb P\) 为有限素数集，并记
\[
G_S:=\ZZ[S^{-1}].
\]
则连续异常相位群满足
\[
\boxed{\;
\mathcal A(G_S)=0.\;}
\]
\end{theorem}
\begin{proof}
任取 \(x,y\in G_S\)。由于 \(x,y\) 都是分母素因子仅属于 \(S\) 的有理数，存在正整数 \(d\)（其全部素因子都属于 \(S\)）使
\[
x,y\in d^{-1}\ZZ.
\]
而
\[
d^{-1}\ZZ\cong \ZZ
\]
是秩 \(1\) 的自由阿贝尔群，故
\[
\wedge^2(d^{-1}\ZZ)=0.
\]
于是 \(x\wedge y=0\) 在 \(\wedge^2G_S\) 中成立。由于 \(\wedge^2G_S\) 由所有简单楔积 \(x\wedge y\) 生成，遂得
\[
\wedge^2G_S=0.
\]
另一方面，由定理 \ref{thm:cdim-anomaly-quadratic-law} 的证明中所调用的经典分解，对交换群与可除系数 \(\TT\) 有
\[
\mathcal A(G)=H^2(G,\TT)\cong \operatorname{Hom}(\wedge^2G,\TT).
\]
代入 \(G=G_S\) 即得
\[
\mathcal A(G_S)\cong \operatorname{Hom}(0,\TT)=0.
\]
证毕。
\end{proof}

\begin{theorem}[有限个局部化数系的异常相位完全分解为 pairwise-union solenoid 之积]\label{thm:conclusion-localized-anomaly-pairwise-union-product}
\leanverified{paper\_conclusion\_localized\_anomaly\_pairwise\_union\_product}
设
\[
G:=\bigoplus_{i=1}^r G_{S_i},
\qquad
G_{S_i}:=\ZZ[S_i^{-1}],
\]
其中每个 \(S_i\subset\mathbb P\) 都是有限素数集。则有自然同构
\[
\boxed{\;
\mathcal A(G)\cong \prod_{1\le i<j\le r}\widehat{G_{S_i\cup S_j}}.\;}
\]
特别地，
\[
\boxed{\;
\cdim_{\QQ}^{\vee}\bigl(\mathcal A(G)\bigr)=\binom r2.\;}
\]
并且对每个 pairwise-union 因子取其自然圆商，可得到规范满射
\[
\boxed{\;
\pi_{\mathrm{Lie}}:\mathcal A(G)\twoheadrightarrow \TT^{\binom r2}.\;}
\]
\end{theorem}
\begin{proof}
仍由上一定理证明末段所用的异常相位分解，
\[
\mathcal A(G)\cong \operatorname{Hom}(\wedge^2G,\TT).
\]
对有限直和应用 exterior square 的标准分解，得
\[
\wedge^2\!\Bigl(\bigoplus_{i=1}^r G_{S_i}\Bigr)
\cong
\bigoplus_{i=1}^r\wedge^2G_{S_i}
\oplus
\bigoplus_{1\le i<j\le r}\bigl(G_{S_i}\otimes_{\ZZ}G_{S_j}\bigr).
\]
由定理 \ref{thm:conclusion-localized-single-axis-anomaly-vanishing} 的证明，\(\wedge^2G_{S_i}=0\) 对每个 \(i\) 都成立，故
\[
\wedge^2G
\cong
\bigoplus_{1\le i<j\le r}\bigl(G_{S_i}\otimes_{\ZZ}G_{S_j}\bigr).
\]
再由局部化的标准张量恒等式，
\[
\ZZ[S_i^{-1}]\otimes_{\ZZ}\ZZ[S_j^{-1}]
\cong
\ZZ[(S_i\cup S_j)^{-1}]
=
G_{S_i\cup S_j},
\]
遂得
\[
\wedge^2G\cong \bigoplus_{1\le i<j\le r}G_{S_i\cup S_j}.
\]
于是
\[
\mathcal A(G)
\cong
\operatorname{Hom}\!\left(\bigoplus_{1\le i<j\le r}G_{S_i\cup S_j},\TT\right)
\cong
\prod_{1\le i<j\le r}\widehat{G_{S_i\cup S_j}},
\]
即得第一条。

由 Pontryagin 对偶，
\[
\widehat{\mathcal A(G)}
\cong
\bigoplus_{1\le i<j\le r}G_{S_i\cup S_j}.
\]
再由定理 \ref{thm:killo-localization-circle-dimension-prime-ledger-splitting}，每个 \(G_{S_i\cup S_j}\) 的有理秩都等于 \(1\)。故
\[
\cdim_{\QQ}^{\vee}\bigl(\mathcal A(G)\bigr)
=
\dim_{\QQ}\!\bigl(\widehat{\mathcal A(G)}\otimes_{\ZZ}\QQ\bigr)
=
\binom r2.
\]

最后，由同一定理，每个因子都满足短正合列
\[
0\longrightarrow \prod_{p\in S_i\cup S_j}\ZZ_p
\longrightarrow
\widehat{G_{S_i\cup S_j}}
\longrightarrow
\TT
\longrightarrow 0.
\]
对所有无序对 \((i,j)\) 的圆商取乘积，即得规范满射
\[
\pi_{\mathrm{Lie}}:\mathcal A(G)\twoheadrightarrow \TT^{\binom r2}.
\]
证毕。
\end{proof}

\begin{corollary}[三寄存器体系中不存在 primitive 三体连续异常]\label{cor:conclusion-localized-three-register-no-primitive-triple-anomaly}
对任意三个有限素数集 \(S_1,S_2,S_3\subset\mathbb P\)，有
\[
\boxed{\;
\mathcal A(G_{S_1}\oplus G_{S_2}\oplus G_{S_3})
\cong
\widehat{G_{S_1\cup S_2}}
\times
\widehat{G_{S_1\cup S_3}}
\times
\widehat{G_{S_2\cup S_3}}.\;}
\]
因而
\[
\boxed{\;
\cdim_{\QQ}^{\vee}\!\bigl(\mathcal A(G_{S_1}\oplus G_{S_2}\oplus G_{S_3})\bigr)=3,\;}
\]
并存在规范满射
\[
\boxed{\;
\pi_{\mathrm{Lie}}:\mathcal A(G_{S_1}\oplus G_{S_2}\oplus G_{S_3})\twoheadrightarrow \TT^3.\;}
\]
特别地，该三寄存器体系中的全部连续异常都由三条 pairwise-union 支撑
\[
S_1\cup S_2,\qquad
S_1\cup S_3,\qquad
S_2\cup S_3
\]
完全决定，不出现依赖于 \(S_1,S_2,S_3\) 同时出现的新的 primitive 三体连续异常因子。
\end{corollary}
\begin{proof}
把定理 \ref{thm:conclusion-localized-anomaly-pairwise-union-product} 代入 \(r=3\) 即得。证毕。
\end{proof}

\begin{theorem}[pairwise-union 异常核的 prime-register 重数公式]\label{thm:conclusion-localized-anomaly-prime-ledger-multiplicity}
\leanverified{paper\_conclusion\_localized\_anomaly\_prime\_ledger\_multiplicity}
沿用定理 \ref{thm:conclusion-localized-anomaly-pairwise-union-product} 的记号。对每个无序对 \((i,j)\)（\(1\le i<j\le r\)），令
\[
\pi_{ij}:\widehat{G_{S_i\cup S_j}}\twoheadrightarrow \TT
\]
为定理 \ref{thm:killo-localization-circle-dimension-prime-ledger-splitting} 中的自然圆商，并记其乘积为
\[
\pi_{\mathrm{Lie}}:\mathcal A(G)\twoheadrightarrow \TT^{\binom r2}.
\]
再对每个素数 \(p\) 定义出现次数
\[
n_p:=\#\{\,i\in\{1,\dots,r\}: p\in S_i\,\}.
\]
则 \(\pi_{\mathrm{Lie}}\) 的零维 prime-register 核满足
\[
\boxed{\;
\ker(\pi_{\mathrm{Lie}})
\cong
\prod_{p\in \bigcup_{i=1}^r S_i}\ZZ_p^{\,m_p},\;}
\]
其中
\[
\boxed{\;
m_p
=
\#\bigl\{(i,j):1\le i<j\le r,\ p\in S_i\cup S_j\bigr\}
=
\binom r2-\binom{r-n_p}{2}.\;}
\]
\end{theorem}
\begin{proof}
由定理 \ref{thm:conclusion-localized-anomaly-pairwise-union-product}，
\[
\mathcal A(G)\cong \prod_{1\le i<j\le r}\widehat{G_{S_i\cup S_j}}.
\]
对每个 pairwise-union 因子应用定理 \ref{thm:killo-localization-circle-dimension-prime-ledger-splitting}，得到
\[
0\longrightarrow \prod_{p\in S_i\cup S_j}\ZZ_p
\longrightarrow
\widehat{G_{S_i\cup S_j}}
\overset{\pi_{ij}}{\longrightarrow}
\TT
\longrightarrow 0.
\]
对全部 \((i,j)\) 取乘积，便得
\[
0\longrightarrow
\prod_{1\le i<j\le r}\prod_{p\in S_i\cup S_j}\ZZ_p
\longrightarrow
\mathcal A(G)
\overset{\pi_{\mathrm{Lie}}}{\longrightarrow}
\TT^{\binom r2}
\longrightarrow 0.
\]
因此
\[
\ker(\pi_{\mathrm{Lie}})
\cong
\prod_{1\le i<j\le r}\prod_{p\in S_i\cup S_j}\ZZ_p.
\]
由于 \(\bigcup_iS_i\) 为有限集，上式可以逐素数重排为
\[
\ker(\pi_{\mathrm{Lie}})
\cong
\prod_{p\in \bigcup_{i=1}^rS_i}\ZZ_p^{\,m_p},
\]
其中 \(m_p\) 恰等于满足 \(p\in S_i\cup S_j\) 的无序对 \((i,j)\) 个数。

固定素数 \(p\)。总无序对数为 \(\binom r2\)；其中不含 \(p\) 的对恰好来自那 \(r-n_p\) 个不含 \(p\) 的指标，因此其个数为
\[
\binom{r-n_p}{2}.
\]
两者相减即得
\[
m_p=\binom r2-\binom{r-n_p}{2}.
\]
证毕。
\end{proof}

\begin{theorem}[pairwise-union 异常核的素数占据数可识别性与终端歧义]\label{thm:conclusion-localized-anomaly-prime-ledger-occupancy-recoverability}
\leanverified{paper\_conclusion\_localized\_anomaly\_prime\_ledger\_occupancy\_recoverability}
沿用定理 \ref{thm:conclusion-localized-anomaly-prime-ledger-multiplicity} 的记号。对固定 \(r\ge 2\)，定义
\[
f_r(n):=\binom r2-\binom{r-n}{2}
\qquad (0\le n\le r).
\]
则对每个素数 \(p\) 的异常核指数
\[
m_p=f_r(n_p)
\]
成立，并且：
\begin{enumerate}
\normalfont
  \item \(f_r\) 在 \(\{0,1,\dots,r-1\}\) 上严格递增；
  \item \(f_r(r-1)=f_r(r)=\binom r2\)；
  \item 因而由 \(m_p\) 可唯一恢复占据数 \(n_p\)，唯一例外是
  \[
  m_p=\binom r2,
  \]
  此时只能断定
  \[
  n_p\in\{r-1,r\}.
  \]
\end{enumerate}
\end{theorem}
\begin{proof}
定理 \ref{thm:conclusion-localized-anomaly-prime-ledger-multiplicity} 已给出
\[
m_p=\binom r2-\binom{r-n_p}{2}=f_r(n_p).
\]
直接计算差分，
\[
f_r(n+1)-f_r(n)
=
\binom{r-n}{2}-\binom{r-n-1}{2}
=
r-n-1.
\]
故当 \(0\le n\le r-2\) 时差分严格为正，从而 \(f_r\) 在 \(\{0,\dots,r-1\}\) 上严格递增。取 \(n=r-1\) 时差分为零，故
\[
f_r(r)=f_r(r-1)=\binom r2.
\]
于是除顶层平台 \(\{r-1,r\}\) 外，\(f_r\) 完全可逆，从而 \(m_p\) 唯一决定 \(n_p\)；而当 \(m_p=\binom r2\) 时仅有两种可能 \(n_p=r-1\) 或 \(n_p=r\)。证毕。
\end{proof}

\begin{corollary}[全覆盖素数与余一缺失素数的异常不可区分性]\label{cor:conclusion-localized-anomaly-prime-ledger-full-vs-almost-full-indistinguishable}
\leanverified{paper\_conclusion\_localized\_anomaly\_prime\_ledger\_full\_vs\_almost\_full\_indistinguishable}
若某素数 \(p\) 满足
\[
n_p=r
\qquad\text{或}\qquad
n_p=r-1,
\]
则其对异常核的贡献完全相同，皆为
\[
\ZZ_p^{\binom r2}.
\]
因此，pairwise-union 异常核不能区分“出现在全部寄存器中”的素数与“恰缺失于一个寄存器”的素数。更一般地，若两组支撑族 \((S_i)\)、\((S_i')\) 只在若干素数上交换了这两种状态，而其余占据数 \(n_p\) 完全相同，则相应零维 prime-register 异常核彼此同构。
\end{corollary}
\begin{proof}
由定理 \ref{thm:conclusion-localized-anomaly-prime-ledger-occupancy-recoverability}，
\[
f_r(r)=f_r(r-1)=\binom r2.
\]
故两种占据状态在异常核中对应同一素数指数 \(\binom r2\)。逐素数比较即可得到所述同构判据。证毕。
\end{proof}

\begin{corollary}[连续异常函子的严格二体化与二次维数增长]\label{cor:conclusion-localized-anomaly-strictly-pairwise-quadratic-growth}
\leanverified{paper\_conclusion\_localized\_anomaly\_strictly\_pairwise\_quadratic\_growth}
设
\[
G=\bigoplus_{i=1}^r G_{S_i},
\qquad r\ge 3.
\]
则连续异常群 \(\mathcal A(G)\) 的全部 primitive 连续因子都严格停留在二体层：不存在任何真正依赖于 \(k\ge 3\) 个寄存器同时参与的 primitive 连续异常因子。与此同时，
\[
\dim\!\bigl(\pi_{\mathrm{Lie}}(\mathcal A(G))\bigr)=\binom r2,
\qquad
\operatorname{rank}_{\QQ}(G)=r.
\]
因此当 \(r\ge 4\) 时，
\[
\dim\!\bigl(\pi_{\mathrm{Lie}}(\mathcal A(G))\bigr)=\binom r2>r,
\]
即连续异常的 Lie 可见维数虽可作二次爆炸，其 primitive 体数却始终停留在二体层。
\end{corollary}
\begin{proof}
由定理 \ref{thm:conclusion-localized-anomaly-pairwise-union-product}，
\[
\mathcal A(G)\cong \prod_{1\le i<j\le r}\widehat{G_{S_i\cup S_j}},
\]
故任一连续异常因子都只能经由某个 pairwise-union 因子出现，不可能需要三个或更多寄存器同时作为 primitive 支撑。又同一定理已给出
\[
\cdim_{\QQ}^{\vee}\bigl(\mathcal A(G)\bigr)=\binom r2,
\]
并构造了规范满射
\[
\pi_{\mathrm{Lie}}:\mathcal A(G)\twoheadrightarrow \TT^{\binom r2},
\]
故其最大紧连通 Lie 商维数正是 \(\binom r2\)。而每个 \(G_{S_i}=\ZZ[S_i^{-1}]\) 的有理秩都等于 \(1\)，故
\[
\operatorname{rank}_{\QQ}(G)=r.
\]
于是当 \(r\ge 4\) 时有 \(\binom r2>r\)。证毕。
\end{proof}

\begin{theorem}[有限素数布尔格上的交并--张量双线性化]\label{thm:conclusion-localized-pairwise-union-wedge-bilinearization}
设 \(S,T\subset\mathbb P\) 为有限素数集。则存在自然同构
\[
\boxed{\;
G_S\otimes_{\ZZ}G_T\cong G_{S\cup T}.\;}
\]
更一般地，对任意有限族 \(S_1,\dots,S_r\subset\mathbb P\)，有
\[
\boxed{\;
\wedge^2\!\Bigl(\bigoplus_{i=1}^r G_{S_i}\Bigr)
\cong
\bigoplus_{1\le i<j\le r}G_{S_i\cup S_j}.\;}
\]
特别地，
\[
\boxed{\;
\wedge^2(G_S\oplus G_T)\cong G_{S\cup T}.\;}
\]
\leanverified{paper\_conclusion\_localized\_pairwise\_union\_wedge\_bilinearization}
\end{theorem}
\begin{proof}
定理 \ref{thm:conclusion-localized-single-axis-anomaly-vanishing} 的证明已经给出
\[
\wedge^2G_{S_i}=0
\qquad(\forall\,i).
\]
另一方面，局部化的标准张量恒等式给出
\[
G_S\otimes_{\ZZ}G_T
=
\ZZ[S^{-1}]\otimes_{\ZZ}\ZZ[T^{-1}]
\cong
\ZZ[(S\cup T)^{-1}]
=
G_{S\cup T}.
\]
对有限直和应用 exterior square 的标准分解，
\[
\wedge^2\!\Bigl(\bigoplus_{i=1}^r G_{S_i}\Bigr)
\cong
\bigoplus_{i=1}^r \wedge^2G_{S_i}
\oplus
\bigoplus_{1\le i<j\le r}(G_{S_i}\otimes_{\ZZ}G_{S_j}),
\]
再代入上面两条恒等式，即得
\[
\wedge^2\!\Bigl(\bigoplus_{i=1}^r G_{S_i}\Bigr)
\cong
\bigoplus_{1\le i<j\le r}G_{S_i\cup S_j}.
\]
当 \(r=2\) 时便得到最后一式。证毕。
\end{proof}

\begin{corollary}[有限局部化异常的二体完备性]\label{cor:conclusion-localized-anomaly-pairwise-completeness}
\leanverified{paper\_conclusion\_localized\_anomaly\_pairwise\_completeness}
设
\[
G:=\bigoplus_{i=1}^r G_{S_i}.
\]
则连续异常相位群满足
\[
\mathcal A(G)\cong \prod_{1\le i<j\le r}\widehat{G_{S_i\cup S_j}}.
\]
并且任意异常元都唯一分解为各 pairwise-union 因子上的二体分量之积。特别地：
\begin{enumerate}
\normalfont
  \item 不存在任何新的 primitive 三体或更高阶连续异常因子；
  \item 两个异常元相等，当且仅当它们在所有 pairwise-union 投影
  \[
  \widehat{G_{S_i\cup S_j}}
  \]
  上逐一相等；
  \item 分辨全部连续异常所需的最小连续证书维数恰为
  \[
  \binom r2.
  \]
\end{enumerate}
\end{corollary}
\begin{proof}
定理 \ref{thm:conclusion-localized-anomaly-pairwise-union-product} 已给出
\[
\mathcal A(G)\cong \prod_{1\le i<j\le r}\widehat{G_{S_i\cup S_j}}
\]
与
\[
\cdim_{\QQ}^{\vee}\bigl(\mathcal A(G)\bigr)=\binom r2.
\]
右端是有限个因子的直积，因此每个元素都有唯一坐标分解；两个元素相等当且仅当全部坐标都相等。由此全部连续异常严格停留在 pairwise 层，不出现新的三体或更高阶 primitive 因子。最后一条即圆维公式的证书解释。证毕。
\end{proof}

\leanverified{paper\_conclusion\_localized\_anomaly\_pairwise\_skeleton\_determines\_spectrum}
\begin{theorem}[有限支持异常谱由 pairwise-union 骨架完全决定]\label{thm:conclusion-localized-anomaly-pairwise-skeleton-determines-spectrum}
设
\[
\mathcal S=\{S_1,\dots,S_r\},
\qquad
\mathcal S'=\{S_1',\dots,S_r'\}
\]
为两组有限素数 cover。若存在排列 \(\sigma\in S_r\) 使得对所有 \(1\le i<j\le r\) 都有
\[
S_i\cup S_j=S_{\sigma(i)}'\cup S_{\sigma(j)}',
\]
则存在自然同构
\[
\boxed{\;
\mathcal A\!\Bigl(\bigoplus_{i=1}^r G_{S_i}\Bigr)
\cong
\mathcal A\!\Bigl(\bigoplus_{i=1}^r G_{S_i'}\Bigr).\;}
\]
因此，有限局部化异常谱的全部连续几何仅由 pairwise-union 骨架决定，而不依赖于单个 \(S_i\) 的更细分解。
\end{theorem}
\begin{proof}
由定理 \ref{thm:conclusion-localized-anomaly-pairwise-union-product}，
\[
\mathcal A\!\Bigl(\bigoplus_{i=1}^r G_{S_i}\Bigr)
\cong
\prod_{1\le i<j\le r}\widehat{G_{S_i\cup S_j}},
\]
\[
\mathcal A\!\Bigl(\bigoplus_{i=1}^r G_{S_i'}\Bigr)
\cong
\prod_{1\le i<j\le r}\widehat{G_{S_i'\cup S_j'}}.
\]
假设给出的排列 \(\sigma\) 使两组 pairwise unions 逐一对应，于是每个因子都与另一侧的同一个局部化对偶群自然同构。对全部无序对取乘积，即得两端异常群自然同构。证毕。
\end{proof}


因此，在有限个局部化 prime-register 轴上，连续异常相位的结构不再只是“存在若干自由度”的粗略描述，而是被完全压缩为一条 pairwise-union 链：单轴项严格消失，\(\wedge^2\)-级别的全部二体异常由逐对并集精确双线性化，任一连续异常元都在 pairwise-union 因子上唯一分解，而整个有限支持异常谱又只依赖于 pairwise-union 骨架本身。再结合前述逐素数二项式重数公式与圆维 \(\binom r2\) 的 sharp 预算，局部化数系的异常相位几何在 rank-\(1\) solenoid 世界中获得了严格的二体化图型分解。

\endinput
